\section{Annexe}

\begin{frame}[fragile]
\frametitle{Annexe 1 : Génération des isoclines et trajectoires}
\begin{mintedbox}[fontsize=\tiny]{python}
import numpy as np
import matplotlib.pyplot as plt

a = 0.9
b = 0.5
epsilon = 0.08
Ip = 1
T = 200
dt = 0.01
t = np.arange(0, T, dt)

v = np.zeros(len(t))
w = np.zeros(len(t))
v[0] = 0
w[0] = 0

for i in range(len(t) - 1):
    dv = v[i] - v[i]**3 / 3 - w[i] + Ip
    dw = epsilon * (v[i] + a - b * w[i])
    v[i + 1] = v[i] + dt * dv
    w[i + 1] = w[i] + dt * dw

v_plot = np.linspace(-2.5, 2.5, 500)
w_null2 = (v_plot + a) / b
w_null_v = v_plot - v_plot**3 / 3 + Ip

plt.figure(figsize=(5, 5))
plt.plot(v_plot, w_null_v, label="dv/dt = 0")
plt.plot(v_plot, w_null2, label="dw/dt = 0")
plt.plot(v, w, alpha=0.6, label="Trajectoire")
plt.xlabel("Potentiel membranaire v")
plt.ylabel("Variable de récupération w")
plt.ylim(-1,3)
plt.legend()
plt.grid()
plt.show()
\end{mintedbox}
\end{frame}

\begin{frame}[fragile]
\frametitle{Annexe 2 : Tracé temporel}
\begin{mintedbox}[fontsize=\tiny]{python}
...
for i in range(len(t) - 1):
    dv = v[i] - v[i]**3 / 3 - w[i] + Ip
    dw = epsilon * (v[i] + a - b * w[i])
    v[i + 1] = v[i] + dt * dv
    w[i + 1] = w[i] + dt * dw
    
plt.figure(figsize=(10, 4))
plt.plot(t, v, label="v(t) : potentiel")
plt.plot(t, w, label="w(t) : récupération")
plt.xlabel("Temps")
plt.ylabel("Variables")
plt.legend()
plt.title("Modèle de FitzHugh–Nagumo")
plt.grid()
plt.show()

\end{mintedbox}
\end{frame}

\begin{frame}[fragile]
\frametitle{Annexe 3 : Comparaison avec modèle linéaire}
\begin{mintedbox}[fontsize=\tiny]{python}
T = 150
dt = 0.01
t = np.arange(0, T, dt)
v = np.zeros(len(t))
w = np.zeros(len(t))
v_lin = np.zeros(len(t))
w_lin = np.zeros(len(t))
v[0] = 0
w[0] = 0
v_lin[0] = 0
w_lin[0] = 0

for i in range(len(t) - 1):
    dv = v[i] - v[i]**3/3 - w[i] + Ip
    dw = epsilon * (v[i] + a - b * w[i])
    v[i+1] = v[i] + dt * dv
    w[i+1] = w[i] + dt * dw
    if v_lin[i] < va: dv_lin = c * v_lin[i] - d - w_lin[i] + Ip
    elif v_lin[i] > vb: dv_lin = c * v_lin[i] + d - w_lin[i] + Ip
    else: dv_lin = (2/3) * v_lin[i] - w_lin[i] + Ip
    dw_lin = epsilon * (v_lin[i] + a - b * w_lin[i])
    v_lin[i+1] = v_lin[i] + dt * dv_lin
    w_lin[i+1] = w_lin[i] + dt * dw_lin

v_plot = np.linspace(-3, 3, 500)
plt.figure(figsize=(6, 6))
plt.plot(t,v, label = "cubique")
plt.plot(t,v_lin, label = "linéaire par morceaux")
plt.xlabel('t')
plt.ylabel('v')
plt.grid()
plt.legend( )
plt.title('Allure de la réponse')
plt.show()
\end{mintedbox}
\end{frame}

\begin{frame}[fragile]
\frametitle{Annexe 4-1: Approximation du Laplacien}
\begin{itemize}

\item{Approximation du laplacien par différences finies centrées à cinq
points}
  \[
  \nabla^2 v_{i,j} \approx \frac{v_{i+1,j} + v_{i-1,j} + v_{i,j+1} + v_{i,j-1} - 4v_{i,j}}{\Delta x^2}.
  \]
\end{itemize}
\begin{mintedbox}[fontsize=\small]{python}
def laplacian(Z):
    return (
        np.roll(Z, 1, axis=0) +
        np.roll(Z, -1, axis=0) +
        np.roll(Z, 1, axis=1) +
        np.roll(Z, -1, axis=1) -
        4 * Z
    ) / dx**2
\end{mintedbox}
\end{frame}

\begin{frame}[fragile]
\frametitle{Annexe 4-2: Initialisation du Tissu et des variables}

\begin{mintedbox}[fontsize=\footnotesize]{python}
import numpy as np
import matplotlib.pyplot as plt

Nx, Ny = 256, 256
dx = 1.0
dt = 0.01
steps = 15000
plot_every = 40

D = 1.0
a = 0.8
b = 0.3
epsilon = 0.08
epsilon_Field = epsilon * (1 + 0.0 * np.random.randn(Nx, Ny))
v = -1.0 * np.ones((Nx, Ny))  
w = -0.5 * np.ones((Nx, Ny))  

v[50:150,Ny//2-4-20:Ny//2+4-10] = 2.0
w[50:150,Ny//2-4-30:Ny//2+4-10] = 0.0
\end{mintedbox}
\end{frame}

\begin{frame}[fragile]
\frametitle{Annexe 4-3: Simulation et visualisation}
\begin{mintedbox}[fontsize=\footnotesize]{python}
plt.figure(figsize=(6, 6))
img = plt.imshow(v, cmap="turbo", origin="lower", animated=True, vmin=-2, vmax=2)
plt.colorbar(label="Potentiel Membranaire")
plt.title("Ondes Spirales")

for n in range(steps):
    Lv = laplacian(v)                     
    dv = v - (v**3) / 3 - w + D * Lv      
    dw = epsilon_Field * (v + a - b * w) 
    v += dt * dv
    w += dt * dw
    if n % plot_every == 1:
        img.set_data(v)          
        plt.draw()
        plt.pause(0.0005)

plt.show()
\end{mintedbox}
\end{frame}